\lampiransection{Template Rekap Data dan Koding Tematik}
\label{app:template-koding-tematik}

Lampiran ini digunakan untuk merekap data hasil wawancara, observasi, dokumentasi, ulasan, dan survei. Setelah observasi GrabFood, kode temuan awal dari Bab IV sudah dapat dicatat sebagai dasar triangulasi, IFAS, EFAS, dan matriks SWOT.

\begin{table}[!htbp]
    \centering
    \caption{Kode Sumber Data Penelitian}
    \label{tab:kode-sumber-data-penelitian}
    \footnotesize
    \setlength{\tabcolsep}{4pt}
    \renewcommand{\arraystretch}{0.95}
    \begin{tabularx}{\textwidth}{|Z{2.2cm}|Y|Y|}
        \toprule
        \textbf{Kode} & \textbf{Sumber Data} & \textbf{Contoh Penggunaan} \\ \hline
        \midrule
        W-PM & Wawancara pemilik atau pengelola & W-PM-01 untuk temuan pertama dari pemilik \\ \hline
        W-AD & Wawancara admin toko & W-AD-01 untuk temuan pertama dari admin \\ \hline
        W-KR & Wawancara karyawan atau dapur & W-KR-01 untuk temuan operasional \\ \hline
        W-PL & Wawancara pelanggan & W-PL-01 untuk temuan dari pelanggan \\ \hline
        W-KP & Wawancara kompetitor atau usaha sejenis & W-KP-PK untuk Pondok Krakatau dan W-KP-IK untuk Istana Krakatau \\ \hline
        W-BC & Wawancara intern \textit{business consultant} Nasi Gerilya & W-BC-01 untuk temuan peran pendamping dan W-BC-21 untuk temuan tambahan informan \\ \hline
        O-GF & Bukti observasi GrabFood Nasi Gerilya & O-GF-20260613 untuk bukti visual observasi Nasi Gerilya \\ \hline
        O-KP & Bukti observasi kompetitor & O-KP-20260619 untuk 209 bukti visual Paripurna, Dendeng Batokok, Istana Krakatau, Pondok Gurih, dan Garuda \\ \hline
        TO-GF & Temuan observasi GrabFood Nasi Gerilya & TO-GF-01 untuk temuan identitas dan reputasi toko \\ \hline
        TO-KP & Temuan observasi kompetitor & TO-KP-01 untuk temuan reputasi kompetitor \\ \hline
        TW-KR/TW-AD & Temuan wawancara pegawai & TW-KR untuk dapur dan TW-AD untuk kasir/checker \\ \hline
        W-BC & Temuan wawancara intern \textit{business consultant} & W-BC untuk sistem antrean, menu digital, stamp digital, dan prioritas operasional \\ \hline
        TW-KP & Temuan wawancara kompetitor & TW-KP-PK dan TW-KP-IK untuk pembanding kompetitor \\ \hline
        TS-PL & Temuan survei/angket pelanggan & TS-PL-01 untuk temuan pelanggan pertama \\ \hline
        D-KIN & Dokumentasi kinerja usaha & D-KIN-01 untuk data penjualan atau indeks \\ \hline
        S-PL & Survei pelanggan pendukung & S-PL-01 untuk ringkasan hasil survei \\ \hline
        \bottomrule
    \end{tabularx}
    \tablesource{Diolah peneliti sebagai sistem kode data penelitian (2026).}
\end{table}

\begin{table}[!htbp]
    \centering
    \caption{Rekap Koding Awal Observasi GrabFood Nasi Gerilya (Bagian 1)}
    \label{tab:rekap-koding-awal-observasi-grabfood}
    \scriptsize
    \setlength{\tabcolsep}{3pt}
    \renewcommand{\arraystretch}{0.92}
    \begin{tabularx}{\textwidth}{|Z{1.6cm}|Y|Z{2.3cm}|Z{2.1cm}|Y|}
        \toprule
        \textbf{Kode Temuan} & \textbf{Ringkasan Temuan} & \textbf{Tema 7P / STP} & \textbf{Kategori Awal} & \textbf{Catatan Analisis} \\ \hline
        \midrule
        TO-GF-01 & Rating 4,7 dari 4.584 penilaian, status toko, kategori, dan jarak 6,6 km terlihat pada profil Nasi Gerilya. & \textit{Physical evidence}, \textit{place} & Kekuatan / catatan akses & Reputasi kuat dari jumlah penilaian, tetapi jarak dan status toko perlu dibaca pada waktu observasi. \\ \hline
        TO-GF-02 & Ringkasan ulasan agregat menonjolkan rasa, porsi, bahan segar, kemasan, variasi menu, dan harga terjangkau. & \textit{Product}, \textit{price} & Kekuatan & Perlu konfirmasi pelanggan karena ringkasan platform dapat menyaring ulasan tertentu. \\ \hline
        TO-GF-03 & Terdapat 19 promo/voucher aktif yang terlihat. & \textit{Promotion}, \textit{price} & Peluang / kekuatan & Promo dapat menjadi pendorong pembelian, tetapi dampaknya pada margin perlu wawancara pemilik. \\ \hline
        TO-GF-04 & Menu terlaris yang terlihat adalah Nasi Padang Rendang Sapi, Nasi Padang Ayam Goreng Bumbu, dan Nasi Padang Ayam Rendang. & \textit{Product}, STP & Kekuatan & Menu terlaris dapat menjadi fokus positioning dan paket promosi. \\ \hline
        TO-GF-05 & Rentang harga keseluruhan Rp1.100--Rp275.000; item tersedia Rp1.100--Rp72.000; paket nasi Padang tersedia Rp27.500--Rp72.000. & \textit{Price} & Kekuatan / risiko & Rentang lebar memberi pilihan, tetapi harga atas harus dibaca bersama ketersediaan dan nilai porsi. \\ \hline
        \bottomrule
    \end{tabularx}
    \tablesource{Diolah peneliti berdasarkan observasi GrabFood tanggal 13 Juni 2026 dan observasi kompetitor tanggal 19 dan 20 Juni 2026.}
\end{table}

\begin{table}[!htbp]
    \centering
    \caption{Rekap Koding Awal Observasi GrabFood Nasi Gerilya dan Kompetitor (Bagian 2)}
    \label{tab:rekap-koding-awal-observasi-grabfood-lanjutan}
    \scriptsize
    \setlength{\tabcolsep}{3pt}
    \renewcommand{\arraystretch}{0.92}
    \begin{tabularx}{\textwidth}{|Z{1.8cm}|Y|Z{2.3cm}|Z{2.1cm}|Y|}
        \toprule
        \textbf{Kode Temuan} & \textbf{Ringkasan Temuan} & \textbf{Tema 7P / STP} & \textbf{Kategori Awal} & \textbf{Catatan Analisis} \\ \hline
        \midrule
        TO-GF-06 & Terdapat 6 menu atau produk yang terlihat tidak tersedia. & \textit{Process}, \textit{product} & Kelemahan & Perlu konfirmasi pengelolaan stok dan pembaruan menu aplikasi. \\ \hline
        TO-GF-07 & Ulasan positif menonjolkan rasa, porsi, kesegaran, kebersihan, pengemasan, request terpenuhi, dan pelanggan rutin. & \textit{Product}, \textit{process} & Kekuatan & Mendukung faktor kekuatan produk dan potensi loyalitas. \\ \hline
        TO-GF-08 & Ulasan negatif/netral memuat salah item, kelengkapan, add-on, persepsi harga/porsi, kualitas, koordinasi staf, dan pengalaman pascakonsumsi. & \textit{Process}, \textit{people}, \textit{product} & Kelemahan & Harus ditriangulasikan dengan admin/karyawan dan pelanggan. \\ \hline
        \begin{tabular}[t]{@{}l@{}}TO-KP-01--\\TO-KP-07\end{tabular} & Lima kompetitor memiliki rating, jumlah penilaian, jarak, promo, rentang harga, tema ulasan, dan penanda reputasi yang menjadi pembanding Nasi Gerilya. & \textit{Price}, \textit{place}, \textit{physical evidence} & Ancaman / peluang & Pondok Gurih unggul reputasi, Garuda dan Istana Krakatau unggul akumulasi penilaian, Dendeng menjadi pembanding harga bawah, dan Paripurna lebih dekat dari titik observasi. \\ \hline
        \bottomrule
    \end{tabularx}
    \tablesource{Diolah peneliti berdasarkan observasi GrabFood tanggal 13 Juni 2026 dan observasi kompetitor tanggal 19 dan 20 Juni 2026.}
\end{table}

\begin{table}[!htbp]
    \centering
    \caption{Rekap Koding Tematik Data Penelitian Utama}
    \label{tab:template-koding-tematik}
    \footnotesize
    \setlength{\tabcolsep}{3pt}
    \renewcommand{\arraystretch}{0.95}
    \begin{tabularx}{\textwidth}{|Z{1.8cm}|Y|Y|Y|Y|}
        \toprule
        \textbf{Kode Data} & \textbf{Ringkasan Temuan} & \textbf{Tema 7P / STP} & \textbf{Kategori SWOT} & \textbf{Catatan Analisis} \\ \hline
        \midrule
        TW-KR-01--TW-KR-10 & Dapur memiliki pembagian kerja, standar resep, cek stok, dan double-check rasa, tetapi rawan keterlambatan saat ramai, menu habis, komunikasi HT tidak jelas, dan keterbatasan SDM. & \textit{Product}, \textit{people}, \textit{process} & S / W & Menjadi dasar faktor S1, S3, W2, W3, dan W5. \\ \hline
        TW-AD-01--TW-AD-12 & Kasir/checker mengelola bill, box, nomor antrean, update lauk, catatan khusus, driver, dan promo; titik rawan muncul pada item terpisah, nomor mirip, dan order membeludak saat promo. & \textit{Place}, \textit{promotion}, \textit{process}, \textit{physical evidence} & S / W / O & Menjadi dasar faktor S3, S4, S5, W1, W4, O2, dan O3. \\ \hline
        W-BC-01--W-BC-21 & Intern \textit{business consultant} menguatkan kebutuhan sistem antrean, menu digital, stamp digital, absensi digital, pengelolaan data pelanggan, paket menu hero, perhitungan promo terhadap HPP/margin, dan indikator evaluasi layanan. & STP, \textit{price}, \textit{place}, \textit{promotion}, \textit{process}, \textit{physical evidence} & S / W / O / T & Menjadi dasar faktor S1, S3, S4, W2, W4, W5, O1, O2, O3, O5, T3, dan T4. \\ \hline
        TW-KP-PK-01--TW-KP-PK-09 & Pondok Krakatau kuat pada rasa dan pelayanan cepat; GrabFood sekitar separuh pesanan, tetapi kompetitor juga mengalami stok belum update, porsi/rasa tidak konsisten, dan item tertinggal. & \textit{Product}, \textit{price}, \textit{place}, \textit{process} & O / T & Menjadi dasar faktor O4, T2, dan T5. \\ \hline
        TW-KP-IK-01--TW-KP-IK-09 & Istana Krakatau kuat pada kanal offline, rasa, service cepat, porsi online, dan sambal; kendala yang muncul adalah pesanan kurang lengkap dan driver menunggu. & \textit{Product}, \textit{price}, \textit{place}, \textit{process} & O / T & Menjadi dasar faktor O4, T1, T2, dan T5. \\ \hline
        TS-PL-01--TS-PL-12 & Pelanggan membeli karena menu unggulan, porsi, rating, foto, jarak, dan promo; keluhan atau risiko ada pada harga, item kurang, stok, suhu, detail menu, dan ketergantungan voucher. & STP, \textit{product}, \textit{price}, \textit{promotion}, \textit{physical evidence} & S / W / O / T & Menjadi dasar faktor S1, S2, W1, W4, W5, O1, O2, O3, T2, T3, dan T4. \\ \hline
        S1--W5 & Faktor internal mencakup rasa/menu unggulan, porsi, SOP online, reputasi digital, respons keluhan, digitalisasi internal, akurasi packing, antrean lintas kanal, stok, harga, informasi menu, dan data pelanggan. & IFAS & S / W & Dipakai pada Matriks IFAS Bab IV. \\ \hline
        O1--T5 & Faktor eksternal mencakup kebutuhan pelanggan GrabFood, promo, tampilan/menu digital, stamp digital, kelemahan kompetitor, media sosial, reputasi kompetitor, pembanding harga, ongkir, biaya platform, bahan baku, dan risiko ulasan negatif. & EFAS & O / T & Dipakai pada Matriks EFAS Bab IV. \\ \hline
        \bottomrule
    \end{tabularx}
    \tablesource{Diolah peneliti berdasarkan transkrip pegawai, wawancara intern \textit{business consultant}, rekap kompetitor, angket pelanggan, observasi GrabFood, dan dokumentasi penelitian utama (2026).}
\end{table}
